%% =====================================================================
%%  preamble.tex  --  how the notes LOOK.   (maths shortcuts: macros.tex)
%%  Assumes  \documentclass[11pt,reqno,oneside]{amsart}  and that main.tex
%%  defines \notetitle and \noteauthor before this file is input.
%% =====================================================================

%% ---- Type ------------------------------------------------------------
\usepackage[T1]{fontenc}
\usepackage[sc]{mathpazo}            % Palatino text + maths, real small caps
\linespread{1.05}                    % Palatino runs tight; this is the usual cure
\usepackage{microtype}               % margin kerning + font expansion

%% ---- Maths -----------------------------------------------------------
\usepackage{amsmath,amssymb,amsthm,mathtools}
\numberwithin{equation}{section}     % (1.3) instead of one global counter

%% ---- Page ------------------------------------------------------------
%% ~75 characters a line (was ~95).  The space this frees on the right is a
%% margin column for \mnote{..} and \todo{..}.  For a plain full-width page:
%%   \usepackage[a4paper, margin=1in, headheight=15pt, footskip=0.4in]{geometry}
\usepackage[a4paper, top=1in, bottom=1in, left=1in, right=1.5in,
            marginparwidth=1.6in, marginparsep=0.3in,
            headheight=15pt, footskip=0.4in]{geometry}

\setlength{\parindent}{0pt}
\setlength{\parskip}{6pt plus 1pt minus 1pt}   % without this, paragraphs run together

%% ---- Colour: one quiet accent, one ink for links ----------------------
\usepackage{xcolor}
\definecolor{accent}{HTML}{8A2F2F}   % oxblood
\definecolor{ink}{HTML}{2C3E6B}      % links
\definecolor{quiet}{gray}{0.40}      % margin notes

\usepackage{graphicx}
\usepackage{enumitem}
\usepackage{tikz}
\usepackage{tikz-cd}
\usetikzlibrary{positioning}
\usepackage{marginnote}
\usepackage{fancyhdr}
% \usepackage{subfiles}

%% ---- Links -----------------------------------------------------------
\usepackage[colorlinks=true, linkcolor=ink, urlcolor=ink, citecolor=ink]{hyperref}
\hypersetup{pdftitle={\notetitle}, pdfauthor={\noteauthor}}
\usepackage{bookmark}
\usepackage[capitalise,noabbrev]{cleveref}   % \cref{thm:x} -> "Theorem 1.2"
\crefformat{equation}{(#2#1#3)}              % \cref{eq:x}  -> "(1.2)"

%% ---- Headings: flush left, numbers hang in the left margin ------------
%% (built on amsart's own hooks; titlesec breaks amsart's contents page)
\makeatletter
\DeclareRobustCommand{\hangnum}[1]{\llap{\textcolor{accent}{#1}\hspace{0.9em}}}
\renewcommand{\@seccntformat}[1]{\hangnum{\csname the#1\endcsname}}
\renewcommand{\section}{\@startsection{section}{1}{\z@}%
  {2.2\linespacing \@plus .6\linespacing \@minus .3\linespacing}{.6\linespacing}%
  {\normalfont\Large\scshape\raggedright}}
\renewcommand{\subsection}{\@startsection{subsection}{2}{\z@}%
  {1.6\linespacing \@plus .4\linespacing \@minus .2\linespacing}{.4\linespacing}%
  {\normalfont\Large\itshape\raggedright}}

%% ---- Title block: flush left, small caps; contents heading to match ---
\renewcommand{\@settitle}{%
  \begin{flushleft}
    \baselineskip14\p@\relax
    \scshape\LARGE\@title
  \end{flushleft}}
\renewcommand{\@setauthors}{%
  \begin{flushleft}
    \itshape\large\authors
  \end{flushleft}}
\renewcommand{\contentsnamefont}{\raggedright\Large\scshape}
\makeatother

%% ---- Running head: title left, current section right; no footer rule ---
\pagestyle{fancy}
\fancyhf{}
%% amsart never calls \sectionmark, so hook it in where it writes the contents entry.
\makeatletter
\let\@amsart@tocwrite\@tocwrite
\renewcommand{\@tocwrite}[2]{%
  \@amsart@tocwrite{#1}{#2}%
  \@ifundefined{#1mark}{}{\csname #1mark\endcsname{#2}}}
\renewcommand{\sectionmark}[1]{%
  \markright{\ifx\@secnumber\@empty\else\ \fi #1}}
\let\subsectionmark\@gobble      % head follows sections only (fancyhdr would also follow subsections)
\makeatother
\fancyhead[L]{\small\scshape\notetitle}
\fancyhead[R]{\small\scshape \rightmark}
\fancyfoot[C]{\small\thepage}
\renewcommand{\headrulewidth}{0.3pt}
\renewcommand{\footrulewidth}{0pt}

%% ---- Contents: no paragraph gaps between entries ----------------------
\makeatletter
\let\@origtableofcontents\tableofcontents
\renewcommand{\tableofcontents}{{\setlength{\parskip}{0pt}\@origtableofcontents}}
\makeatother
